% Generated by code/integrate.py; edit extraction rules there.
% Reviewed source: ECTA.tex, line 64, commit f5021cefa71492babfcbfa580e0c984f59a9de26
\subsection{Relation to the literature}
Neural residual methods build on numerical dynamic programming and function approximation, not on a new principle of optimality. Projection and related methods exploit economic structure to approximate a Bellman equation; see \citet{judd1998}. Derivative approximation by smooth neural networks provides useful representation results \citep{hornik1990}, but does not establish that a particular finite architecture or optimizer solves a nonlinear boundary-value problem.

The differential critic is closely related to the neural PDE approach in \citet{sirignano2018}. Their DGM formulation includes differential and boundary conditions and treats HJB equations. The deep BSDE method of \citet{han2018} also includes an HJB application. The distinguishing object here is the combination of an explicit feasible actor, block-specific evaluation, and economic error accounting across control and equilibrium problems. An analytical formula for the maximizing action is helpful when available, but is not a universal dividing line between neural PDE methods. Likewise, recursive dependence on the value is a nonlinear zeroth-order term; full nonlinearity in the PDE sense concerns the highest derivatives.

Recursive utility separates intertemporal substitution from risk aversion \citep{epsteinzin1989,duffieepstein1992}. Its infinite-horizon formulation requires particular care about utility-process selection and transversality; \citet{herdegen2021} explains why a stationary algebraic solution alone is insufficient. The time-inconsistency extension is an equilibrium problem in the sense of \citet{bjork2016}, not an ordinary control problem with a relabelled flow payoff. Finally, stochastic approximation results require their own step-size, stability, and limiting-dynamics conditions \citep{borkar2024}. They do not convert constant-step Adam into a convergence theorem for economic optimality.

The article proceeds from the economic problem to its evaluation and
improvement map. Section~\ref{sec:theory} states the general verification
results. Sections~\ref{sec:r10continuous}--\ref{sec:r15mechanism} specialize
the welfare account to the capital economy. Section~\ref{sec:economic-scope}
states the other economic applications. The computational study distinguishes
the existing comparison from the current experiment. The supplement collects
the complete models, derivations, proofs, and supporting diagnostics by subject.
